\chapter{Macierze incydencji, cykle i przekroje}
\label{ch:tg-algebra}

W rozdziale \ref{ch:tg-drzewa} dowiedzieliśmy się, że liczba niezależnych
cykli grafu o $n$ wierzchołkach, $m$ krawędziach i $c$ składowych wynosi
$m-n+c$. Po bloku algorytmicznym i spektralnym wyjaśnimy tę formułę
za pomocą macierzy incydencji i algebry liniowej.
Pozwoli to ponadto przejść od zbiorów krawędzi do łańcuchów całkowitych,
co będzie pierwszym rzeczywistym zastosowaniem teorii homologii.

\section{Orientacja pomocnicza i operator brzegowy}
\label{sec:tg-boundary}

W całym rozdziale $G=(V,E)$ oznacza skończony graf prosty
nieskierowany, z $V\ne\varnothing$. Piszemy $n=|V|$, $m=|E|$,
a przez $c$ oznaczamy liczbę składowych spójnych.
\emph{Pomocnicza orientacja} polega na wybraniu dla każdej krawędzi
$e=\{u,v\}$ jednego z dwóch uporządkowań końców: $e:u\to v$ albo
$e:v\to u$. Nie zmieniamy w ten sposób grafu nieskierowanego;
wybór jest jedynie potrzebny do ustalenia znaków.

\begin{definicja}[Łańcuchy wymiaru zero i jeden]
Niech $R$ będzie jednym z pierścieni $\mathbb Z$, $\mathbb Q$ lub
$\mathbb R$. Definiujemy wolne moduły
$C_0(G;R):=R^V$ i $C_1(G;R):=R^E$, których bazy standardowe
oznaczamy odpowiednio przez $[v]$, $v\in V$, oraz $[e]$, $e\in E$.
Dla pomocniczej orientacji $e:u\to v$ określamy $R$-liniowy
\emph{operator brzegowy} $\partial_1:C_1(G;R)\to C_0(G;R)$ przez
$\partial_1[e]=[v]-[u]$.
\end{definicja}

Można myśleć o $[e]$ jako o odcinku posiadającym początek i koniec.
W operatorze brzegowym koniec występuje z plusem, a początek z minusem.
Jeśli odwrócimy pomocniczą orientację pewnych krawędzi,
odpowiadające im kolumny zmienią znak. Niech $D$ będzie
macierzą diagonalną, której współczynnik wynosi $-1$ dla
odwróconej krawędzi i $1$ dla pozostałych. Ponieważ $D^2=I$,
nowa macierz to $BD$, a $BDx=0$ dokładnie wtedy, gdy
$Dx\in\ker B$, czyli $x\in D\ker B$.
Rząd pozostaje ten sam, a łańcuchy zmieniają jedynie
współrzędne względem nowej orientacji.

\begin{definicja}[Macierz incydencji]
Ustalmy kolejność $v_1,\ldots,v_n$ wierzchołków i $e_1,\ldots,e_m$
krawędzi oraz orientację każdej krawędzi. Macierzą incydencji
$B_G=(b_{ij})\in\mathbb Z^{n\times m}$ nazywamy macierz operatora
$\partial_1$, to znaczy $b_{ij}=1$, gdy $v_i$ jest końcem $e_j$,
$b_{ij}=-1$, gdy jest początkiem $e_j$, i $b_{ij}=0$ w pozostałych
przypadkach.
\end{definicja}

\begin{center}
\begin{tikzpicture}[scale=.97,>=Stealth,line cap=round]
 \node[tgBlueVertex] (a) at (-2,0) {$a$};
 \node[tgBlueVertex] (b) at (0,1.4) {$b$};
 \node[tgBlueVertex] (c) at (2,0) {$c$};
 \draw[tgTree,-{Stealth[length=2.8mm]}] (a)--node[above left,font=\small] {$e_1$}(b);
 \draw[tgTree,-{Stealth[length=2.8mm]}] (b)--node[above right,font=\small] {$e_2$}(c);
 \draw[tgAccent,-{Stealth[length=2.8mm]}] (a)--node[below,font=\small] {$e_3$}(c);
 \node[text=tgNavy,font=\footnotesize] at (0,-.92)
   {Wybór strzałek ustala znaki, a nie typ grafu.};
\end{tikzpicture}
\end{center}

\begin{przyklad}[Trójkąt: macierz i brzeg cyklu]
Dla orientacji $e_1:a\to b$, $e_2:b\to c$, $e_3:a\to c$ mamy
\[
B_G=\begin{pmatrix}-1&0&-1\\1&-1&0\\0&1&1\end{pmatrix}.
\]
Suma współrzędnych każdej kolumny wynosi zero. Wektor $z=(1,1,-1)^T$ spełnia
$B_Gz=0$: jest to zamknięty obieg $a\to b\to c\to a$, w którym
trzecią krawędź przechodzimy przeciwnie do jej orientacji pomocniczej.
Natomiast $B_G(1,0,0)^T=(-1,1,0)^T$ opisuje łańcuch mający początek
w $a$ i koniec w $b$.
\end{przyklad}

\section{Rząd macierzy i homologia zerowymiarowa}
\label{sec:tg-rank-incidence}

\begin{twierdzenie}[Rząd macierzy incydencji]
\label{thm:tg-rank}
Niech $G$ będzie skończonym grafem prostym o $n\ge1$ wierzchołkach,
$m$ krawędziach i $c$ składowych spójnych. Niech $B_G$ będzie
macierzą incydencji dla dowolnej pomocniczej orientacji. Wtedy
$\operatorname{rank}_{\mathbb R}B_G=n-c$ i
$\dim_{\mathbb R}\ker B_G=m-n+c$.
\end{twierdzenie}

\begin{proof}
Oznaczmy zbiory wierzchołków składowych przez $V_1,\ldots,V_c$.
Każda kolumna $B_G$ ma sumę współrzędnych równą zero na każdej
składowej. Zatem obraz $B_G$ zawiera się w podprzestrzeni
\[
W=\left\{x\in\mathbb R^V:\sum_{v\in V_i}x_v=0
\text{ dla }i=1,\ldots,c\right\}.
\]
Warunki definiujące $W$ są niezależne, ponieważ zbiory $V_i$ są
niepuste i rozłączne; zatem $\dim W=n-c$.

Pokażemy teraz równość $\operatorname{im}B_G=W$.
W każdej składowej wybierzmy korzeń $r_i$ i drzewo rozpinające.
Dla $v\in V_i\setminus\{r_i\}$ istnieje jednoznaczna ścieżka
w drzewie od $r_i$ do $v$. Sumując jej krawędzie ze współczynnikami
$+1$ albo $-1$ według zgodności z orientacją pomocniczą, otrzymujemy
łańcuch $z_v\in C_1(G;\mathbb R)$ o brzegu
$\partial_1z_v=[v]-[r_i]$: wkłady pośrednich wierzchołków znoszą się
parami. Wektory $[v]-[r_i]$, dla wszystkich $v\ne r_i$, tworzą bazę
$W$: współrzędne poza korzeniami wymuszają niezależność,
a współrzędna przy każdym korzeniu jest ujemną sumą pozostałych
współrzędnych w jego składowej. W szczególności $W\subseteq\operatorname{im}B_G$, co daje
pierwszą tezę. Z twierdzenia o wymiarze jądra i obrazu otrzymujemy
$\dim\ker B_G=m-(n-c)=m-n+c$.
\end{proof}

\begin{uwaga}[Przypadki brzegowe]
Dla pojedynczego wierzchołka bez krawędzi macierz ma wymiar
$1\times0$, rząd zero i jądro zerowe. Dla grafu mającego tylko
wierzchołki izolowane zachodzi $c=n$ i $B_G=0$.
Nie trzeba więc wyłączać tych przypadków z twierdzenia.
\end{uwaga}

\begin{definicja}[Homologie grafu]
Dla skończonego grafu $G$ i $R\in\{\mathbb Z,\mathbb Q,\mathbb R\}$
przyjmujemy $C_k(G;R)=0$ dla $k\notin\{0,1\}$.
Definiujemy $H_1(G;R):=\ker\partial_1$ oraz
$H_0(G;R):=C_0(G;R)/\operatorname{im}\partial_1$.
Są to grupy (odpowiednio moduły) homologii komórkowej
jednowymiarowej realizacji geometrycznej $|G|$.
\end{definicja}

Powyższa definicja jest szczególnym przypadkiem ogólnego kompleksu
łańcuchowego. Ponieważ $C_2=0$, grupa pierwszych brzegów
$\operatorname{im}\partial_2$ jest zerowa, stąd $H_1=\ker\partial_1$.

\begin{twierdzenie}[Homologia grafu ze współczynnikami całkowitymi]
\label{thm:tg-homology}
Niech $G$ będzie skończonym grafem prostym o $n$ wierzchołkach,
$m$ krawędziach i $c$ składowych spójnych. Wówczas
\[
H_0(G;\mathbb Z)\cong\mathbb Z^c,\qquad
H_1(G;\mathbb Z)\cong\mathbb Z^{m-n+c},\qquad
H_k(G;\mathbb Z)=0\quad(k\ge2).
\]
Izomorfizmy nie wymagają wyboru orientacji krawędzi, choć
konkretne bazy pierwszej homologii zależą od wyboru drzewa.
\end{twierdzenie}

\begin{proof}
Rozważmy homologię zerową. Dla każdej składowej $V_i$ określamy
$\varepsilon_i:\mathbb Z^V\to\mathbb Z$, sumując współczynniki
wierzchołków należących do $V_i$. Łączny homomorfizm
$\varepsilon=(\varepsilon_1,\ldots,\varepsilon_c)$ jest surjektywny:
wektor $[r_i]$ dla wybranego korzenia $r_i\in V_i$ przechodzi
na $i$-ty wektor bazy standardowej $\mathbb Z^c$.
Homomorfizm znika na $\operatorname{im}\partial_1$. Ponadto jego jądro jest
równe temu obrazowi: jeśli $x\in\mathbb Z^V$ ma sumę współczynników
zero na każdej składowej, to można zapisać
$x=\sum_i\sum_{v\in V_i\setminus\{r_i\}}x_v([v]-[r_i])$.
Każdy składnik jest brzegiem całkowitego łańcucha odpowiadającego
ścieżce drzewowej od $r_i$ do $v$. Zatem
$\mathbb Z^V/\operatorname{im}\partial_1\cong\mathbb Z^c$.

Dla homologii pierwszej wybieramy las rozpinający $T$.
Dla każdej krawędzi $e\in E\setminus E(T)$ niech $z_e$
będzie całkowitym łańcuchem złożonym z $e$ ze współczynnikiem
$1$ i ze ścieżki w $T$ prowadzącej od końca $e$ do jego początku,
ze współczynnikami $+1$ lub $-1$ zależnie od zgodności przejścia
z orientacją pomocniczą. Taka ścieżka istnieje, ponieważ $T$
rozpina każdą składową spójną. Jeśli $e:u\to v$, jej brzeg
wynosi $[u]-[v]$, a brzeg $[e]$ wynosi $[v]-[u]$.
Wtedy $\partial_1z_e=0$. Wektory $z_e$ są liniowo
niezależne nad $\mathbb Z$, gdyż współrzędna przy krawędzi
$f\notin E(T)$ wektorów $z_e$ wynosi $\delta_{ef}$.

Aby udowodnić, że tworzą bazę całego jądra, weźmy
$z\in\ker\partial_1$ i odejmijmy od niego łańcuch
$\sum_{e\notin E(T)}z_e\,z(e)$, gdzie $z(e)\in\mathbb Z$
jest współczynnikiem $z$ przy $e$. Różnica jest cyklem,
którego nośnik leży w lesie $T$. Taki cykl musi być zerowy:
jeśli jego nośnik jest niepusty, wybieramy liść tego nośnika;
współczynnik przy jego jedynej krawędzi pojawia się w brzegu
z niezerowym znakiem, sprzecznie z $\partial_1z=0$.
Zatem $\{z_e:e\notin E(T)\}$ jest bazą całkowitego
jądra, liczącą $m-(n-c)$ elementów. W wyższych wymiarach
brak komórek, więc pozostałe grupy znikają.
\end{proof}

\section{Przestrzeń cykli modulo dwa i przestrzeń cięć}
\label{sec:tg-cut-cycle}

Do opisu cyklu jako zbioru krawędzi nie potrzebujemy znaków.
Wystarczy przejść do $\mathbb F_2$, w którym $-1=1$.
Niech $B_2$ oznacza macierz incydencji zredukowaną modulo $2$.

\begin{stwierdzenie}[Porównanie z przestrzenią cykli]
Niech $G$ będzie skończonym grafem prostym i niech
$\mathcal C(G)$ będzie przestrzenią cykli zdefiniowaną
w rozdziale \ref{ch:tg-drzewa}. Wtedy
$\mathcal C(G)=\ker(B_2:\mathbb F_2^E\to\mathbb F_2^V)$.
W szczególności ma ona wymiar $m-n+c$.
\end{stwierdzenie}

\begin{proof}
Dla zbioru krawędzi $F\subseteq E$ współrzędna
$B_2\chi_F$ przy wierzchołku $v$ jest równa
stopniowi $v$ w $(V,F)$ modulo $2$.
Każdy cykl prosty ma wszystkie stopnie parzyste, a więc
$\mathcal C(G)\subseteq\ker B_2$.
Odwrotnie, wektor należący do jądra odpowiada zbiorowi
krawędzi o parzystych stopniach. Lemat o rozkładzie takiego
zbioru na cykle, udowodniony w rozdziale
\ref{ch:tg-drzewa}, daje jego przynależność do
$\mathcal C(G)$. Wymiar otrzymujemy z twierdzenia
o bazie cykli fundamentalnych.
\end{proof}

\begin{definicja}[Przestrzeń przekrojów]
Dla grafu $G=(V,E)$ nad $\mathbb F_2$ określamy
$\mathcal D(G)\subseteq\mathbb F_2^E$ jako przestrzeń
liniową rozpiętą przez funkcje charakterystyczne
przekrojów $\delta_G(S)$, gdzie $S\subseteq V$.
\end{definicja}

\begin{twierdzenie}[Ortogonalność cykli i przekrojów]
\label{thm:tg-orthogonal}
Dla dowolnego skończonego grafu prostego $G$ zachodzi
\[
\mathcal D(G)=\operatorname{im}B_2^T
  =\mathcal C(G)^\perp,\qquad
\mathcal C(G)=\mathcal D(G)^\perp,
\]
gdzie ortogonalność jest rozumiana względem
standardowego iloczynu współrzędnych nad $\mathbb F_2$.
Ponadto $\dim\mathcal D(G)=n-c$.
\end{twierdzenie}

\begin{proof}
Dla $S\subseteq V$ wektor $B_2^T\chi_S$
ma przy krawędzi $uv$ współrzędną równą
$\chi_S(u)+\chi_S(v)$ modulo dwa. Jest ona równa
jeden dokładnie wtedy, gdy krawędź przecina podział
$S\sqcup(V\setminus S)$. Każda funkcja charakterystyczna przekroju należy więc do
$\operatorname{im}B_2^T$. Odwrotnie, każdy $f\in\mathbb F_2^V$
jest funkcją charakterystyczną pewnego zbioru $S$.
Stąd każdy element $\operatorname{im}B_2^T$ sam odpowiada
przekrojowi i otrzymujemy $\mathcal D(G)=\operatorname{im}B_2^T$.

Jeśli $z\in\ker B_2$ i $y=B_2^Tf$, to
$\langle z,y\rangle=\langle B_2z,f\rangle=0$.
Stąd $\mathcal D(G)\subseteq\mathcal C(G)^\perp$.
Rząd macierzy incydencji nad $\mathbb F_2$ jest nadal
równy $n-c$: każda kolumna ma sumę zero na każdej składowej,
a wektory $[v]+[r_i]$ dla $v\in V_i\setminus\{r_i\}$
są obrazami ścieżek i pozostają liniowo niezależne,
co daje oba oszacowania rzędu.
Tak więc $\dim\mathcal D(G)=n-c$, podczas gdy
$\dim\mathcal C(G)^\perp=m-(m-n+c)=n-c$.
Daje to pierwszą równość, a drugą otrzymujemy przez
ponowne przejście do dopełnienia ortogonalnego,
ponieważ standardowe parowanie na $\mathbb F_2^E$
jest niezdegenerowane.
\end{proof}

\begin{przyklad}[Kwadrat z przekątną: dwa niezależne cykle]
Niech $V=\{1,2,3,4\}$, a krawędzie będą skierowane
$e_1:1\to2$, $e_2:2\to3$, $e_3:3\to4$,
$e_4:4\to1$, $e_5:1\to3$. Wówczas
\[
B=\begin{pmatrix}
-1&0&0&1&-1\\1&-1&0&0&0\\
0&1&-1&0&1\\0&0&1&-1&0
\end{pmatrix}.
\]
Dla drzewa rozpinającego $T=\{e_1,e_2,e_3\}$
dwa fundamentalne cykle mają wektory
$z_4=(1,1,1,1,0)^T$ i $z_5=(-1,-1,0,0,1)^T$.
Bezpośrednie mnożenie daje $Bz_4=Bz_5=0$.
Ponieważ $m-n+1=5-4+1=2$, tworzą one bazę
$H_1(G;\mathbb Z)\cong\mathbb Z^2$.

W wersji modulo dwa mamy wektory
$\bar z_4=(1,1,1,1,0)$ oraz $\bar z_5=(1,1,0,0,1)$.
Ich suma odpowiada trójkątowi na wierzchołkach $1,3,4$.
Wyróżnione dwa cykle są bazą, ale nie są wszystkimi
cyklami grafu.
\end{przyklad}

\begin{center}
\begin{tikzpicture}[scale=1.0,line cap=round]
 \node[tgBlueVertex] (a) at (-1.6,1.0) {$1$};
 \node[tgBlueVertex] (b) at (1.6,1.0) {$2$};
 \node[tgBlueVertex] (c) at (1.6,-1.0) {$3$};
 \node[tgBlueVertex] (d) at (-1.6,-1.0) {$4$};
 \draw[tgTree] (a)--node[above,font=\scriptsize] {$e_1$}(b);
 \draw[tgTree] (b)--node[right,font=\scriptsize] {$e_2$}(c);
 \draw[tgTree] (c)--node[below,font=\scriptsize] {$e_3$}(d);
 \draw[tgAccent] (d)--node[left,font=\scriptsize] {$e_4$}(a);
 \draw[tgAccent] (a)--node[pos=.68,above,font=\scriptsize] {$e_5$}(c);
 \node[text=tgNavy,font=\footnotesize] at (0,-1.53)
   {Zielone: drzewo $T$. Czerwone: krawędzie generujące cykle.};
\end{tikzpicture}
\end{center}

\begin{uwaga}[Ortogonalność nad $\mathbb F_2$ nie oznacza sumy prostej]
Równości $\mathcal C(G)=\mathcal D(G)^\perp$ oraz
$\mathcal D(G)=\mathcal C(G)^\perp$ nie implikują,
że ich przecięcie jest zerowe. Przykładowo w cyklu
$C_4$ zbiór wszystkich czterech krawędzi jest cyklem,
a jednocześnie jest przekrojem $\delta(\{v_1,v_3\})$.
Zatem należy odróżnić te równości nad $\mathbb F_2$
od rzeczywistego ortogonalnego rozkładu na gradienty i cykle.
\end{uwaga}

\section{Całkowite macierze incydencji i przepływy}
\label{sec:tg-incidence-integral}

\begin{definicja}[Całkowita unimodularność]
Macierz całkowitą nazywamy całkowicie unimodularną, jeżeli
wyznacznik każdej jej kwadratowej podmacierzy należy do
$\{-1,0,1\}$. Wyznacznik pustej macierzy przyjmujemy równy $1$.
\end{definicja}

\begin{twierdzenie}[Całkowita unimodularność macierzy incydencji]
\label{thm:tg-incidence-tu}
Macierz incydencji dowolnego skończonego grafu prostego,
dla każdej pomocniczej orientacji, jest całkowicie unimodularna.
\end{twierdzenie}
\begin{proof}
Niech $M$ będzie kwadratową podmacierzą macierzy incydencji.
Stosujemy indukcję względem jej rozmiaru $k$, zaczynając od $k=0$.
Jeśli pewna kolumna jest zerowa, to $\det M=0$.
Jeśli pewna kolumna ma dokładnie jeden niezerowy element,
jest on równy $1$ lub $-1$. Rozwinięcie Laplace'a względem
tej kolumny sprowadza wyznacznik, z dokładnością do znaku,
do wyznacznika podmacierzy rozmiaru $k-1$.
W pozostałym przypadku każda kolumna zawiera dokładnie dwa
niezerowe elementy, jeden równy $1$, drugi równy $-1$.
Suma wszystkich wierszy jest wtedy zerowa, a ponieważ $k>0$,
wiersze są liniowo zależne i $\det M=0$.
\end{proof}

\begin{wniosek}[Całkowitość wierzchołków wielościanu przepływów]
\label{cor:tg-integral-flow-vertices}
Niech $B$ będzie macierzą incydencji skończonego grafu prostego,
$b\in\mathbb Z^V$ oraz $\ell,u\in\mathbb Z^E$,
przy czym $\ell_e\le u_e$ dla każdego $e\in E$.
Każdy punkt ekstremalny zbioru
$P=\{x\in\mathbb R^E:Bx=b,\ \ell\le x\le u\}$
ma całkowite współrzędne. Nierówności są rozumiane współrzędnie.
\end{wniosek}
\begin{proof}
Niech $x$ będzie punktem ekstremalnym $P$ i niech
$F=\{e:\ell_e<x_e<u_e\}$. Poza $F$ współrzędne $x$ są
równe całkowitym ograniczeniom. Kolumny $B$ odpowiadające
$F$ są liniowo niezależne: w przeciwnym razie istniałby
niezerowy wektor $h$ o nośniku zawartym w $F$ i $Bh=0$.
Ze skończoności $F$ i ścisłych nierówności na $F$ wynika,
że dla dostatecznie małego $\varepsilon>0$ oba punkty
$x+\varepsilon h$ i $x-\varepsilon h$ należą do $P$.
Są różne, a ich średnią jest $x$, wbrew ekstremalności.

Jeżeli $F$ jest pusty, teza już zachodzi. W przeciwnym razie
wybieramy $|F|$ wierszy, na których kolumny $F$ tworzą
odwracalną macierz kwadratową $M$. Taki wybór istnieje
z równości rzędu wierszowego i kolumnowego.
Po przeniesieniu znanych współrzędnych na prawą stronę
otrzymujemy układ $Mx_F=d$ z całkowitym $d$.
Na mocy twierdzenia~\ref{thm:tg-incidence-tu} mamy
$\det M=\pm1$. Wzór Cramera daje całkowite współrzędne $x_F$.
\end{proof}

\begin{przyklad}[Ograniczenia przepływu na trójkącie]
W orientacji z pierwszego przykładu przyjmijmy
$b=(-1,0,1)^T$ oraz $0\le x_e\le1$.
Równania $Bx=b$ dają $x_1=x_2=t$, $x_3=1-t$,
gdzie $0\le t\le1$. Zbiór dopuszczalny jest odcinkiem
o końcach $(0,0,1)$ i $(1,1,0)$.
Oba końce są całkowite, choć punkt środkowy nie jest.
Wniosek~\ref{cor:tg-integral-flow-vertices} dotyczy
punktów ekstremalnych, a nie wszystkich przepływów.
\end{przyklad}

\section{Rozkład ortogonalny w przestrzeni krawędzi}
\label{sec:tg-linear-hodge}

Zastąpmy teraz $\mathbb F_2$ przez $\mathbb R$ i nadajmy
przestrzeniom $C_0(G;\mathbb R)$, $C_1(G;\mathbb R)$
standardowe iloczyny skalarne, w których bazy wierzchołków
i krawędzi są ortonormalne. Wtedy operatorem sprzężonym
do $B$ jest $B^T$.

\begin{twierdzenie}[Ortogonalny rozkład Hodge'a na grafie]
\label{thm:tg-edge-hodge}
Dla skończonego grafu prostego $G$ zachodzi suma prosta ortogonalna
\[
\mathbb R^E=\operatorname{im}B^T\ \oplus^\perp\ \ker B.
\]
W szczególności każdy łańcuch krawędziowy $x$ można
zapisać jako $x=B^Tf+h$ z $Bh=0$.
Składniki $B^Tf$ i $h$ są wyznaczone jednoznacznie.
Sam wektor $f$ jest wyznaczony z dokładnością do dodania funkcji
stałej na każdej składowej grafu.
\end{twierdzenie}

\begin{proof}
Dla $h\in\ker B$ i dowolnego $f\in\mathbb R^V$ mamy
$\langle h,B^Tf\rangle=\langle Bh,f\rangle=0$.
Zatem obie podprzestrzenie są ortogonalne.
Z twierdzenia \ref{thm:tg-rank} wynika, że
$\dim\ker B=m-n+c$, natomiast
$\dim\operatorname{im}B^T=\operatorname{rank}B^T
=\operatorname{rank}B=n-c$.
Suma tych wymiarów jest równa $m$,
więc otrzymujemy rozkład całej przestrzeni.

Z rozkładu wynika istnienie $f$ oraz jednoznaczność
składników $B^Tf$ i $h$. Równanie $BB^Tf=Bx$
jest równoważne warunkowi $B(x-B^Tf)=0$.
Ponadto $\ker(BB^T)=\ker B^T$, gdyż
$\langle BB^Tf,f\rangle=\|B^Tf\|^2$.
Stąd $\operatorname{rank}(BB^T)=\operatorname{rank}B$
i $\operatorname{im}(BB^T)=\operatorname{im}B$,
gdyż pierwszy obraz zawiera się w drugim, a oba mają
ten sam wymiar.
Dwa rozwiązania $f,g$ różnią się o element $\ker B^T$,
czyli funkcję stałą na każdej składowej (jej wartości
na końcach każdej krawędzi są równe).
\end{proof}

Przestrzenie cykli nad $\mathbb F_2$ i $\mathbb R$
mają ten sam wymiar, lecz różnią się współczynnikami:
nad $\mathbb R$ orientacja ustala znaki, a nad $\mathbb F_2$
znaki znikają. Jądro operatora brzegowego opisuje $H_1$.

